\chapter{对Hall结果的判决及下一步研究}
\section{模型能力比较：新增自由度、实际成果与代价}
\small
\begin{longtable}{P{26mm}P{39mm}P{43mm}P{34mm}}
\toprule 模型 & 新增可计算量 & 已有直接证据 & 仍受限制\\\midrule\endfirsthead
\toprule 模型 & 新增可计算量 & 已有直接证据 & 仍受限制\\\midrule\endhead
共速双温 & 两个温度、交换滞后、传导前沿 & L164的碰撞激波时空结构 & 给定出流，不求扩散区；经典离子闭合可失效\\
双温耗散Hall & 独立物种预算、壁面热损与电子对流 & L223、HIT-SI的温图和趋势 & 热源分配/输运系数/边界；非完整张量\\
惯性五矩 & 电子流能、惯性响应、两套标量总能 & Wang2015及MMS五矩对照 & 无各向异性；MMS例能量损失较大\\
局地十矩 & 六分量压力、非对角支撑、各向异性 & Wang2015；TenBarge2019MMS全域热能 & 零迹松弛难替代真实热输运；峰侧别错\\
非局地十矩 & 近似尺度依赖相混及压力输运 & Ng2017离子层与压力支撑改善 & 无磁化闭合边界；电子结构仍偏差\\
梯度十矩 & 热流散度、岛反弹、部分横向模态 & L256、Ng2020、E552、Closure2026 & 参数校准、跨场扩散、正定与三维收敛\\
各向异性EoS & 给定来源下平行/垂直电子温度 & Wetherton2019/2021观测检验 & 未提供通用离子闭合；局地俘获历史遗漏\\
Landau/陀螺流体 & 部分平行相混、惯性与有限回旋半径 & L244的增长率/细结构变化 & 实际冷离子；热化延迟与自由能归宿\\
来源分群电子 & 群间漂移对总压力和热流的贡献 & F635电子独有重联示例 & 固定离子；群标签、晚期多束及幅度偏差\\
十三/二十矩 & 理论上演化热流并闭合四阶矩 & 本轮只追到待核查线索 & 不据矩数或会议摘要声称加热验证\\
\bottomrule\end{longtable}
\normalsize

\section{三个比较问题的逐项判决}
\subsection{能得到Hall得不到的结论吗}
\textbf{相对H0/H1：能。}独立能量使温度比例和滞后不再由输入规定；L164与HIT-SI实际展示了重要差别。新增方程是两套热力学，而不是Hall电磁项。

\textbf{相对标量H2/H3：部分能。}十矩能描述非对角压力、各向异性温度和相应应变功；EoS能预测由俘获限制热流所造成的温度--密度转折。这些是标量压力体系无法直接输出的量。其正确性仍取决于闭合，尤其是热流。模型能表示某一量与模型已经准确预测该量必须分开。

\textbf{相对H4：没有由名称保证的优势。}如果Hall约化已保留相同压力张量、热流和独立能量，区别主要转到电子惯性、电荷分离、位移电流及高频尺度。在低频准中性适用区，完整Maxwell多流体应逼近相应Hall极限。大幅偏差首先要求核查参数、边界、能量一致性与数值误差，而非自动视为新物理。

\subsection{能和实验对上吗}
可以对上部分温度形态、趋势和电子闭合关系；本轮证据没有支持“完整双流体已同时精确匹配实验$T_i,T_e$、张量压力、热流、两物种不可逆加热和总能预算”。实验验证最强的部分集中在有限控制体预算或特定电子状态方程，而不是一套跨装置闭合。原文未给统一误差指标，不能以本报告的表格整齐代替这一缺口。

\subsection{能独立印证Hall加热结果吗}
支持程度必须按独立性排序：同一模型的两份代码可检查实现；不同模型共用同一热源分配和调参，只能提供有限的稳健性证据；真正不同闭合在未用于校准的观测量上同时成功，才对物理解释有更强支持。L164与MAST下游离子先热的路径相容，但碰撞性、几何和供能方式不同，不能当作同参数重复验证。十矩闭合的离子层变化与MMS各向异性偏差则限制把Hall峰温相符等同于正确机制的做法。

\section{已解决与未解决：四类缺口分别处理}
\begin{longtable}{P{40mm}P{56mm}P{52mm}}
\toprule 问题 & 已有求解尝试及条件 & 剩余缺口与类别\\\midrule\endfirsthead
\toprule 问题 & 已有求解尝试及条件 & 剩余缺口与类别\\\midrule\endhead
温差怎样产生与消失 & 两温能量、碰撞交换、热传导；L164/HIT-SI & 快交换与无碰撞过渡；输运闭合可信度\\
离子出流怎样变成压力 & 十矩离子压力支撑、岛合并反弹；Ng/L256 & 缺实验压力功和热流联合验证；证据不足\\
电子为何平行更热 & 俘获EoS，强导场MMS & 弱导场非旋转对称压力与局地历史；模型边界\\
三维波动是否改变分配 & E552/Closure2026恢复部分模态 & 未给收敛物种熵产；目标尚未检验\\
多束与热核如何分开 & F635按电子来源分群 & 离子、三维及动态分群未验证；表示限制\\
Hall能否得到通用加热比 & 装置双温Hall与实验预算 & 源、边界、几何依赖且观测定义不同；无此证据\\
更多矩是否更好 & 原理上可演化热流 & 十三/二十矩完整证据未取得；文献覆盖缺口\\
原始各向异性流体链 & Ohia2012/2015原始论文已定位 & 正文/图式访问缺口未闭合，不据摘要入结论\\
\bottomrule\end{longtable}

前六项是本报告依据跨文献作出的综合判断，部分与原文明确限制重合；最后两项是本轮证据覆盖的缺口。未取得原文不能写成学界没有做过，模型不能表示尾部也不意味着所有加热问题都需要全动理学。

\section{研究设计一：先把模型差异归因到一条新增方程}
选择两个可控问题：无导场Fadeev岛合并用于离子压力支撑；强导场近对称重联用于电子温度各向异性。第一阶段不做装置全尺度，优先复用已公开输入，逐项核对归一化、边界和扰动，而非只抄论文表格。

同一初态与同一边界下依次计算：H1单热力学、H2独立能量、H3电子惯性、五矩全Maxwell、局地十矩、梯度十矩；强导场支线再与EoS比较。H1需明确指定分配比例，不能把该比例与其他模型的输出混合统计。H3与全Maxwell的差异要单独隔离；若同时换代码，应先做共享物理极限验证。

\textbf{消融顺序：}固定电子惯性，改变标量/张量压力；固定压力模型，改变热流；固定闭合参数，改变质量比和光速；最后才改变碰撞、导场和几何。碰撞交换、黏性及电阻必须在能量中成对记账。对局地松弛，额外输出它对压力迹的零贡献与对矩熵的贡献，避免把两者混同。

\textbf{输出：}各物种$W_s,\Delta K_s,\Delta U_s$、压缩功、无迹应变功、边界焓流、热流、碰撞交换、物理耗散、数值残差；$T_{\parallel},T_\perp$、压力特征值、热点位置与宽度。分别按相同时间和相同重联磁通比较，以区分演化时序与状态结构。

\textbf{成败门槛（本报告建议，非文献已有结果）：}三档网格及至少两档时间步，在最细两档上主要积分预算和热点宽度变化小于5\%；全能量残差与分物种预算残差各小于同期释放磁能的1\%，并小于所讨论模型差值的五分之一；压力保持正定，限制器/下限带来的能量改变量单列。若差异小于数值不确定度，不报告“新增物理改变加热”。这些门槛应在试算前固定，必要时因实验精度调整并留记录。

\section{研究设计二：不用同一条温度曲线同时调参与验证}
选择MAST的一组可获取放电或MMS一个公开事件。MAST需在模拟中实现实验驱动和热损边界，合成碳离子谱与电子散射响应；MMS需沿候选轨迹采样$n,B,T_{e\parallel},T_{e\perp}$，保留航天器时间到空间映射的不确定性。原始数据未取得前不得承诺已可重现。

先划分校准与验证对象：只用上游$n,B,T_s$及几何选择边界；若必须用一个温度剖面拟合$k_0$、黏性或来源参数，另一个时刻、位置、导场工况或热流曲线必须留出。固定参数后再预测测试集，不允许逐事件继续调参后声称通用。

对Wetherton2021类三来源拟合，先用入流及独立磁场资料限制来源，盲预测出流温度及紫色偏差区。若只有逐区域自由选择四个参数才能吻合，应报告拟合解释能力；若加入来源记忆变量后能预测未参与拟合的垂直增温，才支持模型扩展的预测价值。

\textbf{成败判据：}模拟与观测在预先指定区间的差值是否落在测量、参数传播与数值误差合成区间内；热点侧别和温度转折是否正确；能量残差是否明显小于待解释热量。原文缺少可用协方差时先报告分项误差，不机械计算约化$\chi^2$。若温度命中而功率/通量失配，判为机制尚未通过。

\section{研究设计三：区分净增温、散射与不可逆性}
对十矩模型，在总温度匹配后，进一步检查压力特征值、各向异性、热流与熵闭合。对正定压力张量，局地Gaussian矩熵每粒子可写为
\begin{equation}
\sigma_G/k_B=\frac12\ln\!\left(\frac{\det\vct P}{n^5}\right)+\text{含单位基准的常数}.
\end{equation}
这是用有限矩重建的Gaussian熵，不是未知真实分布的精细熵。各向同性化在固定迹下提高行列式，可提高$\sigma_G$却不增加内能；此诊断能揭示“散射增加”与“加热增加”的混淆，但不能替代分布测量。

对来源分群模型，同时报告群内内能和群间漂移贡献，检查总矩合成前后的能量恒等式。若温升主要来自群间漂移，论文和报告都应明确为有效二阶矩升高；需要实验可分辨的多束或碰撞松弛证据，才能进一步论证热化。此阶段只在关键点使用少量动理学参照，不扩展成本轮研究主线。

\section{最小可复查交付清单}
每个算例保存初边值、单位、质量比、$c/V_A$、网格、步长、离散格式、闭合参数及限制器，列出物理参数与校准参数。运行记录须保留能量残差、压力最小特征值、源项积分及边界通量。最终结论按“新增方程改变了哪个可测量、幅度超过何种误差、在哪些参数内成立”撰写；失败工况与成功工况一起保留。此报告完成的是文献和研究设计，没有声称执行以上模拟。
